\begin{resumo}
 A crescente digitalização dos serviços de saúde e a sobrecarga dos sistemas de atendimento têm impulsionado a busca por soluções automatizadas para triagem e anamnese clínica. No entanto, a aplicação direta de Modelos de Linguagem de Grande Escala (LLMs) neste domínio enfrenta desafios críticos relacionados à alucinação factual e à falta de estruturação dos dados gerados. Este trabalho propõe e avalia tecnicamente uma arquitetura de software baseada no padrão Backend for Frontend (BFF) para a realização de anamnese assistida. A solução adota uma estratégia de Inteligência Artificial híbrida, orquestrando o modelo generativo Llama-3.1 para a condução empática do diálogo e o modelo extrativo Biomedical-NER para a identificação e categorização de entidades clínicas, como sintomas e medicamentos. A metodologia envolveu levantamento bibliográfico, definição de requisitos, prototipação incremental e implementação de uma API em Node.js integrada a uma interface web reativa. Os resultados preliminares do protótipo funcional indicam viabilidade técnica da abordagem, evidenciando que a separação de responsabilidades entre geração de texto e extração de dados reduz inconsistências e produz saídas estruturadas para integração posterior com prontuários eletrônicos. O trabalho estabelece as bases arquiteturais para a continuação do projeto com expansão funcional, plano de testes mais amplo e validação com usuários.

 \vspace{\onelineskip}
    
 \noindent
 \textbf{Palavras-chave}: Inteligência Artificial na Saúde. Anamnese Assistida. Engenharia de Software. LLMs. NER.
\end{resumo}
